\subsection{1653—1672：海上退守与皇权重组}

福建的港湾没有因为府县已换了旗号便归入同一种秩序。顺治十年（1653年）以后，郑成功以厦门、金门为泊船、转运和集结的支点，在漳州、泉州、同安之间反复争夺水陆相接的地方。清廷在南方取得城池和官署，首先意味着可以设置守备、征催钱粮和接收投诚；它并不使沿海的商人、船户、耕作者和旧有军政关系同时改变。郑氏所能倚赖的，也不是一条抽象的“海上生命线”，而是港口外可停泊的水面、通往腹地的渡口、愿意或被迫供给的村落，以及仍可传递消息的亲族和贸易关系。陆上的城防若能截住这些接点，岛屿便难以养兵；郑军若能保持它们，清方的府县接收便始终有边界。

顺治十六年（1659年）的北上南京，把这种未竟的南方接收推到长江下游。郑军试图把海上的机动与江南的明遗民联系、粮赋资源结合起来，却须在内河水道、城外营地和城防之间维持船只、粮食与联络。撤退并非只是一场战役的结果：它显示一支能够穿入长江的舰队，仍难以在远离岛屿据点的地方承受围城和久驻的耗费。南京周围的行动不能用作江南人心的总表；有人提供消息、有人避战、有人被征发，也有人首先面对的是保城、保粮或逃离的选择。清方战报、杨英《从征实录》、地方记事和荷兰方面的海上消息能够彼此校正北伐与撤军的大致次序，却不足以把内应的范围、兵员损失或每座村镇的立场写成整齐数字。

因此，1661年的迁界令与郑成功登陆台湾必须放在同一张水陆网络中理解。清廷命闽粤居民内迁，所要切断的是郑氏从岸上获得粮食、人口和情报的可能；这是一种以迁徙制造军事距离的省级治理，并不是海面从此封闭的证明。谕旨、督抚文书和地方志可看见命令的层层下达、若干撤迁与后来的复业线索，却不能由此推算整段海岸的迁徙距离、死亡人数或秘密贸易规模。渔盐、田地、渡海、家族和市集原本相互牵连，禁令把风险分配得并不均匀：一处港口的沉寂不能代表另一处，留在原地、迁往内陆和改换生计的人，也不会在官府档案中留下同样清楚的痕迹。

台湾在此时不是南京撤军后自然出现的“退路”，而是争夺粮食、人口和航路的另一处战场。郑军围攻热兰遮城，至1662年初荷兰守军投降；城堡、贸易节点和周边土地随之易手。围城使荷兰东印度公司的记录保存了守军补给、城内决策和交接的一个视角，投降条约也能钉住转换的时点；它们都不能代替城外汉人居民、被俘者或原住民社群的叙述。郑成功旋即去世，继承之争又使新基地必须重新安置军政、宗族和海商网络。1663年厦门、金门易手，郑氏主力转守台湾，海峡两边的压力并未消失，只是更多地落到岛内土地、港口与跨海经营所能供给的限度上。

北京的权力重组同样要穿过这些地方性的供给和文书链。1661年顺治帝去世，八岁的玄烨即位，索尼、苏克萨哈、遏必隆、鳌拜辅政；1665年前后鳌拜影响扩大，1667年索尼去世，原有的均衡更难维持。辅政者处置的并非只是一座宫廷中的职位，还包括旗务、人事、军费以及各省报上来的战情和钱粮。1669年康熙拘捕鳌拜，扩大了皇帝直接统摄军政人事的条件，却不能倒过来证明此后每一项调拨都畅通无阻。1669年案件的案卷和后出的传记尤其容易把结局写成早已注定的个人胜负；满汉题本与起居注留下的议事程序，反而提醒我们当时的约束并未因后来处分而消失。

1670年颁行十六条圣谕，朝廷尝试把伦理、赋役和治安的语言交给地方官绅与社区。这与沿海隔离和战时转运一样，显示中央试图以可传达的规范接续有限的官僚、驻军和财力；但诏令的颁布只能证明国家说了什么，不能证明某县怎样宣讲、何人听见，或社会秩序因而如何改变。各省方志、基层诉讼和契约偶尔保存这类接触的局部痕迹，保存的又多是有文书能力的官、绅与产权关系，不能替无名的役夫、流徙者和船户发言。到1672年，南方的接收、海峡的战争和省级的筹饷仍未收束为稳定秩序；它们留下的驻军、道路、债务与地方中介关系，也构成次年撤藩冲突得以扩大的条件。

\paragraph{材料位置。} 《康熙起居注》、满汉题本和鳌拜案档案呈现议事、处分与命令的文书链；《清史稿》的后出人物叙事则带有1669年以后胜利的视角。沿海一侧，实录和军报须与郑氏文书、地方志、投降条约及荷兰东印度公司热兰遮城日志并读：它们能够连接政策、围城和港口，却各自遮蔽不同社群。中央重组与海岸战争共享财政、交通和信息空间，仍不能彼此化约。

\section{摄政的终结与皇权的重新组织}

古北口外的行猎行程在顺治七年十二月突然中断：多尔衮去世，原先围绕摄政王组织的军政枢纽随之改组（EQ-1650-DORGON-DEATH）。《清世祖实录》和内国史院的丧仪、遗命档可以确认死期和制度性的后果；死因以及死后清算的动机，则不能照抄后来人物传记的定性。次年顺治帝亲政，撤去摄政安排、重置皇帝周边的任官和决策网络（EQ-1651-SHUNZHI-PERSONAL-RULE）。亲政诏令和题本首先证明命令的发出和人事的程序，并不证明南方将领、降官、旗务与部院已经无条件服从。皇帝名义的强化仍须经过同一套军饷、道路和文书链才能抵达地方。

顺治十八年，八岁的玄烨即位，索尼、苏克萨哈、遏必隆、鳌拜辅政（EQ-1661-KANGXI-ACCESSION）。继位和四人名单可由《清世祖实录》、起居注、内国史院遗诏档交叉确立；顺治死因与遗诏形成过程仍有材料和后世传说的层次差异。康熙四年辅政大臣依议政与旗务网络处置军政，鳌拜的影响扩大（EQ-1665-REGENTS-CONSOLIDATION）；索尼于1667年去世，使原有均衡更加脆弱（EQ-1667-SONI-DEATH）。《康熙起居注》与满汉题本呈现的是当时的议事和程序，而《清史稿·鳌拜传》及鳌拜案档案带着1669年胜利后的视角。把后来被拘捕的结局提前写成“独揽大权”，会抹掉辅政本来具有的集体性和相互制约。

康熙八年五月拘捕鳌拜并处分其党羽，的确改变了皇帝直接统摄人事与军政的条件（EQ-1669-OBOI-ARREST）。拘捕、审理和处分可由实录、起居注及内阁案卷相互验证；宫中擒捕的戏剧性细节、罪状的完整性和少年皇帝的单独谋略，则都是胜利叙事最容易放大的部分。康熙九年颁布十六条圣谕（EQ-1670-SACRED-EDICT），又把中央的伦理、赋役和秩序语言交给地方官绅与社区。圣谕的文本和宣讲制度是可见的，识读范围、宣讲频率与社会效果却必须由各省方志、基层文书和具体时段检验。它说明国家试图以训导补足武力和官僚覆盖的不足，而不是一条可以直接量得的服从曲线。


\subsection{事件链与材料限度}
摄政、宗室、部院与地方官之间的文书和人事安排保留了不同速度。

\paragraph{顺治七年十二月（1650年）：摄政王多尔衮在古北口外行猎途中去世，摄政政治结构随之改组}
\eventanchor{EQ-1650-DORGON-DEATH}{摄政王多尔衮在古北口外行猎途中去世，摄政政治结构随之改组}。本节点叙述该事件本身。清在京师与多地军政线路相接的尺度的处境首先受多尔衮集中军政与皇室事务，围绕摄政权的宗室和旗人竞争在其死后失去调停中心制约。《清世祖实录》顺治七年十二月；《清史稿·多尔衮传》；《清初内国史院满文档案》摄政王丧仪与遗命档留下可核的线索；可辨的后果是顺治帝及近臣得以重组朝廷，摄政集团的名誉与权力被重新界定。原有置信注记：死亡日期可靠；死因、身后清算与权力斗争不能只依后来的政治定性

\paragraph{顺治七年十二月（1650年）：摄政王多尔衮在古北口外行猎途中去世，摄政政治结构随之改组}
\eventanchor{EQ-1650-DORGON-DEATH-AFTERMATH}{摄政王多尔衮在古北口外行猎途中去世，摄政政治结构随之改组} 置于京师与多地军政线路相接的尺度中阅读，本节点追踪「摄政王多尔衮在古北口外行猎途中去世，摄政政治结构随之改组」之后可见的余波。多尔衮集中军政与皇室事务，围绕摄政权的宗室和旗人竞争在其死后失去调停中心构成行动者清必须面对的条件。取证从《清世祖实录》顺治七年十二月；《清史稿·多尔衮传》；《清初内国史院满文档案》摄政王丧仪与遗命档出发，因而只能把变化谨慎地连到顺治帝及近臣得以重组朝廷，摄政集团的名誉与权力被重新界定。原有置信注记：死亡日期可靠；死因、身后清算与权力斗争不能只依后来的政治定性

\paragraph{顺治七年十二月（1650年）：摄政王多尔衮在古北口外行猎途中去世，摄政政治结构随之改组}
\eventanchor{EQ-1650-DORGON-DEATH-DECISION}{摄政王多尔衮在古北口外行猎途中去世，摄政政治结构随之改组} 标记本节点定位围绕「摄政王多尔衮在古北口外行猎途中去世，摄政政治结构随之改组」形成的处置与调度记录，涉及清。京师与多地军政线路相接的尺度不是抽象背景：多尔衮集中军政与皇室事务，围绕摄政权的宗室和旗人竞争在其死后失去调停中心。《清世祖实录》顺治七年十二月；《清史稿·多尔衮传》；《清初内国史院满文档案》摄政王丧仪与遗命档所见的顺序随后指向顺治帝及近臣得以重组朝廷，摄政集团的名誉与权力被重新界定，但原有置信注记：死亡日期可靠；死因、身后清算与权力斗争不能只依后来的政治定性

\paragraph{顺治七年十二月（1650年）：摄政王多尔衮在古北口外行猎途中去世，摄政政治结构随之改组}
在京师与多地军政线路相接的尺度，\eventanchor{EQ-1650-DORGON-DEATH-PRELUDE}{摄政王多尔衮在古北口外行猎途中去世，摄政政治结构随之改组} 让清的一个选择或压力有了可查的坐标。本节点回看「摄政王多尔衮在古北口外行猎途中去世，摄政政治结构随之改组」之前已可见的条件。前因可写为多尔衮集中军政与皇室事务，围绕摄政权的宗室和旗人竞争在其死后失去调停中心；《清世祖实录》顺治七年十二月；《清史稿·多尔衮传》；《清初内国史院满文档案》摄政王丧仪与遗命档限定了记录，后续变化是顺治帝及近臣得以重组朝廷，摄政集团的名誉与权力被重新界定。原有置信注记：死亡日期可靠；死因、身后清算与权力斗争不能只依后来的政治定性

\paragraph{顺治八年（1651年）：顺治帝亲政，撤除摄政安排并重置皇帝周边的决策网络}
\eventanchor{EQ-1651-SHUNZHI-PERSONAL-RULE}{顺治帝亲政，撤除摄政安排并重置皇帝周边的决策网络} 所关联的是清在京师与多地军政线路相接的尺度中的行动。本节点叙述该事件本身。它从多尔衮去世使摄政权威无法延续，皇帝及支持者以礼制和任免重建指挥链展开，借《清世祖实录》顺治八年正月；《清史稿·世祖本纪》；《顺治朝内阁档》亲政诏令与任官题本进入可检验的年序，并以中央命令署名和人事结构改变，但南方征服与财政仍要求依靠地方将领和降官影响下一段安排。原有置信注记：亲政事实可靠；实际决策仍受宗室、旗务和满汉官僚共同制约

\paragraph{顺治八年（1651年）：顺治帝亲政，撤除摄政安排并重置皇帝周边的决策网络}
京师与多地军政线路相接的尺度中的\eventanchor{EQ-1651-SHUNZHI-PERSONAL-RULE-AFTERMATH}{顺治帝亲政，撤除摄政安排并重置皇帝周边的决策网络} 不能离开清来解释。本节点追踪「顺治帝亲政，撤除摄政安排并重置皇帝周边的决策网络」之后可见的余波。多尔衮去世使摄政权威无法延续，皇帝及支持者以礼制和任免重建指挥链说明其形成的条件；《清世祖实录》顺治八年正月；《清史稿·世祖本纪》；《顺治朝内阁档》亲政诏令与任官题本提供来源路径，中央命令署名和人事结构改变，但南方征服与财政仍要求依靠地方将领和降官则是材料能够支持的近时结果。原有置信注记：亲政事实可靠；实际决策仍受宗室、旗务和满汉官僚共同制约

\paragraph{顺治八年（1651年）：顺治帝亲政，撤除摄政安排并重置皇帝周边的决策网络}
\eventanchor{EQ-1651-SHUNZHI-PERSONAL-RULE-DECISION}{顺治帝亲政，撤除摄政安排并重置皇帝周边的决策网络} 把清放回京师与多地军政线路相接的尺度的道路、城镇、边界或制度网络。本节点定位围绕「顺治帝亲政，撤除摄政安排并重置皇帝周边的决策网络」形成的处置与调度记录。多尔衮去世使摄政权威无法延续，皇帝及支持者以礼制和任免重建指挥链。本段采用《清世祖实录》顺治八年正月；《清史稿·世祖本纪》；《顺治朝内阁档》亲政诏令与任官题本核对其顺序，随后可追到中央命令署名和人事结构改变，但南方征服与财政仍要求依靠地方将领和降官。原有置信注记：亲政事实可靠；实际决策仍受宗室、旗务和满汉官僚共同制约

\paragraph{顺治八年（1651年）：顺治帝亲政，撤除摄政安排并重置皇帝周边的决策网络}
本节点回看「顺治帝亲政，撤除摄政安排并重置皇帝周边的决策网络」之前已可见的条件，故\eventanchor{EQ-1651-SHUNZHI-PERSONAL-RULE-PRELUDE}{顺治帝亲政，撤除摄政安排并重置皇帝周边的决策网络} 不只是一个年份标签。清在京师与多地军政线路相接的尺度承受的压力来自多尔衮去世使摄政权威无法延续，皇帝及支持者以礼制和任免重建指挥链。《清世祖实录》顺治八年正月；《清史稿·世祖本纪》；《顺治朝内阁档》亲政诏令与任官题本可回查该节点；中央命令署名和人事结构改变，但南方征服与财政仍要求依靠地方将领和降官是其进入后续年序的方式。原有置信注记：亲政事实可靠；实际决策仍受宗室、旗务和满汉官僚共同制约

\paragraph{顺治十八年正月（1661年）：顺治帝去世，八岁玄烨即位为康熙帝，索尼、苏克萨哈、遏必隆、鳌拜辅政}
\eventanchor{EQ-1661-KANGXI-ACCESSION}{顺治帝去世，八岁玄烨即位为康熙帝，索尼、苏克萨哈、遏必隆、鳌拜辅政} 所见的行动者为清，空间尺度为京师与多地军政线路相接的尺度。本节点叙述该事件本身。顺治早逝使未成年皇帝继承，宗室与旗人高层需建立可承认的集体决策安排使它成为具体事件链的一环；《清世祖实录》顺治十八年正月；《清史稿·圣祖本纪》；《康熙起居注》与内国史院遗诏档与辅政集团掌握朝政，皇权成长、旗人权力平衡和南方战事资源相互交织。原有置信注记：继位及辅政名单可靠；顺治死因和遗诏形成过程存在材料与后世叙事差异

\paragraph{顺治十八年正月（1661年）：顺治帝去世，八岁玄烨即位为康熙帝，索尼、苏克萨哈、遏必隆、鳌拜辅政}
从京师与多地军政线路相接的尺度看，\eventanchor{EQ-1661-KANGXI-ACCESSION-AFTERMATH}{顺治帝去世，八岁玄烨即位为康熙帝，索尼、苏克萨哈、遏必隆、鳌拜辅政} 留下本节点追踪「顺治帝去世，八岁玄烨即位为康熙帝，索尼、苏克萨哈、遏必隆、鳌拜辅政」之后可见的余波的材料坐标。清所面对的条件是顺治早逝使未成年皇帝继承，宗室与旗人高层需建立可承认的集体决策安排。《清世祖实录》顺治十八年正月；《清史稿·圣祖本纪》；《康熙起居注》与内国史院遗诏档支持的结果为辅政集团掌握朝政，皇权成长、旗人权力平衡和南方战事资源相互交织。原有置信注记：继位及辅政名单可靠；顺治死因和遗诏形成过程存在材料与后世叙事差异

\paragraph{顺治十八年正月（1661年）：顺治帝去世，八岁玄烨即位为康熙帝，索尼、苏克萨哈、遏必隆、鳌拜辅政}
\eventanchor{EQ-1661-KANGXI-ACCESSION-DECISION}{顺治帝去世，八岁玄烨即位为康熙帝，索尼、苏克萨哈、遏必隆、鳌拜辅政} 记录清在京师与多地军政线路相接的尺度的一次变化。本节点定位围绕「顺治帝去世，八岁玄烨即位为康熙帝，索尼、苏克萨哈、遏必隆、鳌拜辅政」形成的处置与调度记录。其发生与顺治早逝使未成年皇帝继承，宗室与旗人高层需建立可承认的集体决策安排相连；《清世祖实录》顺治十八年正月；《清史稿·圣祖本纪》；《康熙起居注》与内国史院遗诏档把事件系回来源，辅政集团掌握朝政，皇权成长、旗人权力平衡和南方战事资源相互交织显示压力如何向后转移。原有置信注记：继位及辅政名单可靠；顺治死因和遗诏形成过程存在材料与后世叙事差异

\paragraph{顺治十八年正月（1661年）：顺治帝去世，八岁玄烨即位为康熙帝，索尼、苏克萨哈、遏必隆、鳌拜辅政}
\eventanchor{EQ-1661-KANGXI-ACCESSION-PRELUDE}{顺治帝去世，八岁玄烨即位为康熙帝，索尼、苏克萨哈、遏必隆、鳌拜辅政} 的地点与行动范围属于京师与多地军政线路相接的尺度，行动者是清。本节点回看「顺治帝去世，八岁玄烨即位为康熙帝，索尼、苏克萨哈、遏必隆、鳌拜辅政」之前已可见的条件。顺治早逝使未成年皇帝继承，宗室与旗人高层需建立可承认的集体决策安排。读《清世祖实录》顺治十八年正月；《清史稿·圣祖本纪》；《康熙起居注》与内国史院遗诏档时可见其记录边界，最小后果只能写作辅政集团掌握朝政，皇权成长、旗人权力平衡和南方战事资源相互交织。原有置信注记：继位及辅政名单可靠；顺治死因和遗诏形成过程存在材料与后世叙事差异

\paragraph{康熙四年（1665年）：四辅政大臣以议政和旗务网络处理军政事务，鳌拜在政务中的影响扩大}
\eventanchor{EQ-1665-REGENTS-CONSOLIDATION}{四辅政大臣以议政和旗务网络处理军政事务，鳌拜在政务中的影响扩大} 是清在京师与多地军政线路相接的尺度留下的编年标记。本节点叙述该事件本身。它从幼帝尚未亲裁，辅政大臣需在旗人贵族、部院官僚和南方军费之间协调走来；《清圣祖实录》康熙四年政务条；《清史稿·鳌拜传》；《康熙起居注》及内阁大库满汉题本固定可证部分，中央决策呈现集体辅政与个人网络并存，为1669年权力重组提供背景构成下一步的可见结果。原有置信注记：辅政架构可靠；权力高下多由后来的鳌拜案档与《清史稿》倒叙，须避免预设结局

\paragraph{康熙四年（1665年）：四辅政大臣以议政和旗务网络处理军政事务，鳌拜在政务中的影响扩大}
本节点追踪「四辅政大臣以议政和旗务网络处理军政事务，鳌拜在政务中的影响扩大」之后可见的余波：\eventanchor{EQ-1665-REGENTS-CONSOLIDATION-AFTERMATH}{四辅政大臣以议政和旗务网络处理军政事务，鳌拜在政务中的影响扩大}。清在京师与多地军政线路相接的尺度的处境可由幼帝尚未亲裁，辅政大臣需在旗人贵族、部院官僚和南方军费之间协调说明。《清圣祖实录》康熙四年政务条；《清史稿·鳌拜传》；《康熙起居注》及内阁大库满汉题本是本段材料入口，因而只能据此追踪到中央决策呈现集体辅政与个人网络并存，为1669年权力重组提供背景。原有置信注记：辅政架构可靠；权力高下多由后来的鳌拜案档与《清史稿》倒叙，须避免预设结局

\paragraph{康熙四年（1665年）：四辅政大臣以议政和旗务网络处理军政事务，鳌拜在政务中的影响扩大}
\eventanchor{EQ-1665-REGENTS-CONSOLIDATION-DECISION}{四辅政大臣以议政和旗务网络处理军政事务，鳌拜在政务中的影响扩大} 发生在京师与多地军政线路相接的尺度，牵动清。本节点定位围绕「四辅政大臣以议政和旗务网络处理军政事务，鳌拜在政务中的影响扩大」形成的处置与调度记录。幼帝尚未亲裁，辅政大臣需在旗人贵族、部院官僚和南方军费之间协调是其结构性条件；《清圣祖实录》康熙四年政务条；《清史稿·鳌拜传》；《康熙起居注》及内阁大库满汉题本留下的证据把读者带到中央决策呈现集体辅政与个人网络并存，为1669年权力重组提供背景。原有置信注记：辅政架构可靠；权力高下多由后来的鳌拜案档与《清史稿》倒叙，须避免预设结局

\paragraph{康熙四年（1665年）：四辅政大臣以议政和旗务网络处理军政事务，鳌拜在政务中的影响扩大}
\eventanchor{EQ-1665-REGENTS-CONSOLIDATION-PRELUDE}{四辅政大臣以议政和旗务网络处理军政事务，鳌拜在政务中的影响扩大} 对应本节点回看「四辅政大臣以议政和旗务网络处理军政事务，鳌拜在政务中的影响扩大」之前已可见的条件。在京师与多地军政线路相接的尺度，清无法绕开幼帝尚未亲裁，辅政大臣需在旗人贵族、部院官僚和南方军费之间协调。依据《清圣祖实录》康熙四年政务条；《清史稿·鳌拜传》；《康熙起居注》及内阁大库满汉题本，本段把后续影响限定为中央决策呈现集体辅政与个人网络并存，为1669年权力重组提供背景。原有置信注记：辅政架构可靠；权力高下多由后来的鳌拜案档与《清史稿》倒叙，须避免预设结局

\paragraph{康熙六年（1667年）：首辅索尼去世，四辅政大臣的均衡被打破}
\eventanchor{EQ-1667-SONI-DEATH}{首辅索尼去世，四辅政大臣的均衡被打破}。本节点叙述该事件本身。清在京师与多地军政线路相接的尺度的处境首先受索尼居于辅政资历与旗人关系枢纽，其去世改变了原有制衡条件制约。《清圣祖实录》康熙六年；《清史稿·索尼传》；《康熙起居注》康熙六年丧仪与议政记录留下可核的线索；可辨的后果是鳌拜及盟友影响更显著，康熙帝近侍和太皇太后方面开始寻找收回裁决权的机会。原有置信注记：索尼死亡及官职变动可靠；其政治立场与遗言多经后出人物传记修饰

\paragraph{康熙六年（1667年）：首辅索尼去世，四辅政大臣的均衡被打破}
\eventanchor{EQ-1667-SONI-DEATH-AFTERMATH}{首辅索尼去世，四辅政大臣的均衡被打破} 置于京师与多地军政线路相接的尺度中阅读，本节点追踪「首辅索尼去世，四辅政大臣的均衡被打破」之后可见的余波。索尼居于辅政资历与旗人关系枢纽，其去世改变了原有制衡条件构成行动者清必须面对的条件。取证从《清圣祖实录》康熙六年；《清史稿·索尼传》；《康熙起居注》康熙六年丧仪与议政记录出发，因而只能把变化谨慎地连到鳌拜及盟友影响更显著，康熙帝近侍和太皇太后方面开始寻找收回裁决权的机会。原有置信注记：索尼死亡及官职变动可靠；其政治立场与遗言多经后出人物传记修饰

\paragraph{康熙六年（1667年）：首辅索尼去世，四辅政大臣的均衡被打破}
\eventanchor{EQ-1667-SONI-DEATH-DECISION}{首辅索尼去世，四辅政大臣的均衡被打破} 标记本节点定位围绕「首辅索尼去世，四辅政大臣的均衡被打破」形成的处置与调度记录，涉及清。京师与多地军政线路相接的尺度不是抽象背景：索尼居于辅政资历与旗人关系枢纽，其去世改变了原有制衡条件。《清圣祖实录》康熙六年；《清史稿·索尼传》；《康熙起居注》康熙六年丧仪与议政记录所见的顺序随后指向鳌拜及盟友影响更显著，康熙帝近侍和太皇太后方面开始寻找收回裁决权的机会，但原有置信注记：索尼死亡及官职变动可靠；其政治立场与遗言多经后出人物传记修饰

\paragraph{康熙六年（1667年）：首辅索尼去世，四辅政大臣的均衡被打破}
在京师与多地军政线路相接的尺度，\eventanchor{EQ-1667-SONI-DEATH-PRELUDE}{首辅索尼去世，四辅政大臣的均衡被打破} 让清的一个选择或压力有了可查的坐标。本节点回看「首辅索尼去世，四辅政大臣的均衡被打破」之前已可见的条件。前因可写为索尼居于辅政资历与旗人关系枢纽，其去世改变了原有制衡条件；《清圣祖实录》康熙六年；《清史稿·索尼传》；《康熙起居注》康熙六年丧仪与议政记录限定了记录，后续变化是鳌拜及盟友影响更显著，康熙帝近侍和太皇太后方面开始寻找收回裁决权的机会。原有置信注记：索尼死亡及官职变动可靠；其政治立场与遗言多经后出人物传记修饰

\paragraph{康熙八年五月（1669年）：康熙帝拘捕鳌拜并审理其党羽，结束辅政大臣实际主政的局面}
\eventanchor{EQ-1669-OBOI-ARREST}{康熙帝拘捕鳌拜并审理其党羽，结束辅政大臣实际主政的局面} 所关联的是清在京师与多地军政线路相接的尺度中的行动。本节点叙述该事件本身。它从索尼死后辅政均衡失去支点，鳌拜与苏克萨哈案件激化皇帝近臣、太皇太后和旗人贵族冲突展开，借《清圣祖实录》康熙八年五月；《清史稿·鳌拜传》；《康熙起居注》与鳌拜案内阁档进入可检验的年序，并以皇帝可直接统摄人事军政，但仍须通过议政、内阁及旗务结构行使权力影响下一段安排。原有置信注记：拘捕与处分可靠；宫中擒捕过程、罪状与少年皇帝个人谋略含胜利者叙事

\paragraph{康熙八年五月（1669年）：康熙帝拘捕鳌拜并审理其党羽，结束辅政大臣实际主政的局面}
京师与多地军政线路相接的尺度中的\eventanchor{EQ-1669-OBOI-ARREST-AFTERMATH}{康熙帝拘捕鳌拜并审理其党羽，结束辅政大臣实际主政的局面} 不能离开清来解释。本节点追踪「康熙帝拘捕鳌拜并审理其党羽，结束辅政大臣实际主政的局面」之后可见的余波。索尼死后辅政均衡失去支点，鳌拜与苏克萨哈案件激化皇帝近臣、太皇太后和旗人贵族冲突说明其形成的条件；《清圣祖实录》康熙八年五月；《清史稿·鳌拜传》；《康熙起居注》与鳌拜案内阁档提供来源路径，皇帝可直接统摄人事军政，但仍须通过议政、内阁及旗务结构行使权力则是材料能够支持的近时结果。原有置信注记：拘捕与处分可靠；宫中擒捕过程、罪状与少年皇帝个人谋略含胜利者叙事

\paragraph{康熙八年五月（1669年）：康熙帝拘捕鳌拜并审理其党羽，结束辅政大臣实际主政的局面}
\eventanchor{EQ-1669-OBOI-ARREST-DECISION}{康熙帝拘捕鳌拜并审理其党羽，结束辅政大臣实际主政的局面} 把清放回京师与多地军政线路相接的尺度的道路、城镇、边界或制度网络。本节点定位围绕「康熙帝拘捕鳌拜并审理其党羽，结束辅政大臣实际主政的局面」形成的处置与调度记录。索尼死后辅政均衡失去支点，鳌拜与苏克萨哈案件激化皇帝近臣、太皇太后和旗人贵族冲突。本段采用《清圣祖实录》康熙八年五月；《清史稿·鳌拜传》；《康熙起居注》与鳌拜案内阁档核对其顺序，随后可追到皇帝可直接统摄人事军政，但仍须通过议政、内阁及旗务结构行使权力。原有置信注记：拘捕与处分可靠；宫中擒捕过程、罪状与少年皇帝个人谋略含胜利者叙事

\paragraph{康熙八年五月（1669年）：康熙帝拘捕鳌拜并审理其党羽，结束辅政大臣实际主政的局面}
本节点回看「康熙帝拘捕鳌拜并审理其党羽，结束辅政大臣实际主政的局面」之前已可见的条件，故\eventanchor{EQ-1669-OBOI-ARREST-PRELUDE}{康熙帝拘捕鳌拜并审理其党羽，结束辅政大臣实际主政的局面} 不只是一个年份标签。清在京师与多地军政线路相接的尺度承受的压力来自索尼死后辅政均衡失去支点，鳌拜与苏克萨哈案件激化皇帝近臣、太皇太后和旗人贵族冲突。《清圣祖实录》康熙八年五月；《清史稿·鳌拜传》；《康熙起居注》与鳌拜案内阁档可回查该节点；皇帝可直接统摄人事军政，但仍须通过议政、内阁及旗务结构行使权力是其进入后续年序的方式。原有置信注记：拘捕与处分可靠；宫中擒捕过程、罪状与少年皇帝个人谋略含胜利者叙事

\paragraph{康熙九年（1670年）：康熙帝颁布十六条圣谕，规定由地方官和民众遵循的伦理、赋役与秩序训导框架}
\eventanchor{EQ-1670-SACRED-EDICT}{康熙帝颁布十六条圣谕，规定由地方官和民众遵循的伦理、赋役与秩序训导框架} 所见的行动者为清，空间尺度为省、府县与地方社会相互牵动的尺度。本节点叙述该事件本身。清廷在征服后寻求以共同伦理语言补足武力和官僚的有限覆盖，地方官绅被要求传达使它成为具体事件链的一环；《清圣祖实录》康熙九年十月；《清史稿·圣祖本纪》；《大清会典事例》圣谕条与各省地方志宣讲记录与圣谕成为后续地方教化与政治合法性资源，其接受和改写具有地区差异。原有置信注记：颁布可靠；宣讲频率、识读范围和实际社会效果不能由文本证明

\paragraph{康熙九年（1670年）：康熙帝颁布十六条圣谕，规定由地方官和民众遵循的伦理、赋役与秩序训导框架}
从省、府县与地方社会相互牵动的尺度看，\eventanchor{EQ-1670-SACRED-EDICT-AFTERMATH}{康熙帝颁布十六条圣谕，规定由地方官和民众遵循的伦理、赋役与秩序训导框架} 留下本节点追踪「康熙帝颁布十六条圣谕，规定由地方官和民众遵循的伦理、赋役与秩序训导框架」之后可见的余波的材料坐标。清所面对的条件是清廷在征服后寻求以共同伦理语言补足武力和官僚的有限覆盖，地方官绅被要求传达。《清圣祖实录》康熙九年十月；《清史稿·圣祖本纪》；《大清会典事例》圣谕条与各省地方志宣讲记录支持的结果为圣谕成为后续地方教化与政治合法性资源，其接受和改写具有地区差异。原有置信注记：颁布可靠；宣讲频率、识读范围和实际社会效果不能由文本证明

\paragraph{康熙九年（1670年）：康熙帝颁布十六条圣谕，规定由地方官和民众遵循的伦理、赋役与秩序训导框架}
\eventanchor{EQ-1670-SACRED-EDICT-DECISION}{康熙帝颁布十六条圣谕，规定由地方官和民众遵循的伦理、赋役与秩序训导框架} 记录清在省、府县与地方社会相互牵动的尺度的一次变化。本节点定位围绕「康熙帝颁布十六条圣谕，规定由地方官和民众遵循的伦理、赋役与秩序训导框架」形成的处置与调度记录。其发生与清廷在征服后寻求以共同伦理语言补足武力和官僚的有限覆盖，地方官绅被要求传达相连；《清圣祖实录》康熙九年十月；《清史稿·圣祖本纪》；《大清会典事例》圣谕条与各省地方志宣讲记录把事件系回来源，圣谕成为后续地方教化与政治合法性资源，其接受和改写具有地区差异显示压力如何向后转移。原有置信注记：颁布可靠；宣讲频率、识读范围和实际社会效果不能由文本证明

\paragraph{康熙九年（1670年）：康熙帝颁布十六条圣谕，规定由地方官和民众遵循的伦理、赋役与秩序训导框架}
\eventanchor{EQ-1670-SACRED-EDICT-PRELUDE}{康熙帝颁布十六条圣谕，规定由地方官和民众遵循的伦理、赋役与秩序训导框架} 的地点与行动范围属于省、府县与地方社会相互牵动的尺度，行动者是清。本节点回看「康熙帝颁布十六条圣谕，规定由地方官和民众遵循的伦理、赋役与秩序训导框架」之前已可见的条件。清廷在征服后寻求以共同伦理语言补足武力和官僚的有限覆盖，地方官绅被要求传达。读《清圣祖实录》康熙九年十月；《清史稿·圣祖本纪》；《大清会典事例》圣谕条与各省地方志宣讲记录时可见其记录边界，最小后果只能写作圣谕成为后续地方教化与政治合法性资源，其接受和改写具有地区差异。原有置信注记：颁布可靠；宣讲频率、识读范围和实际社会效果不能由文本证明

\section{海岸不是陆上征服的边缘}

厦门与金门之间的水道并不因为陆上府县易手而安静下来。顺治十年（1653年），郑成功以两岛为支点，把军队、商船和仍可调动的沿海关系拢在同一套经营中。岛屿提供停泊与转运的地点，港口把粮食、人员和消息送向舰队；但要把这些东西送到前线，还得越过潮汐、水道与清军据守的城镇。郑氏因此不能只守海面：漳州、泉州、同安一线的攻守，关系到腹地粮食能否抵达海边，也关系到城防能否把这条补给线截断。

顺治十一年至十三年间，围攻漳州以及在泉州、同安附近争夺水陆通道，正是这种压力的展开。郑军需要扩大粮饷与招募的基础，清军则依托府县城防和陆路援兵，使其难以迅速取得大陆府城。对港口周边的人们而言，战争不是两支舰队在远海相遇：征粮、避乱、治安的变化，会落在村落、渡口、盐场和市集不同的位置。我们可以从军报、地方志与郑氏文书看见攻守和若干地点，却很难据此排出每一处居民如何选择、何时离开或继续交易的整齐次序。商人、船民、耕作者、守城者和逃亡者面对的风险不必相同，更不能被概括成一种一致的“支持”或“抵抗”。

这也是郑氏海上力量的两面。商船和侨民网络能够输送筹饷的机会，岛屿也能让军政机构避开陆上包围；可是船只、兵员和收入的数字多见于敌对军报或后出的郑氏叙事，不能把它们当作精确的家底。较为稳妥的判断是：福建海疆已无法仅由陆上州县来控制。清廷后来把海禁、迁界与水师建设连在一起，并不是出于抽象的禁海观念，而是要处理港口、岛屿、腹地与消息网络彼此相连的难题。

顺治十六年（1659年）北上南京，使这种难题沿江而上。郑成功试图利用长江下游的粮赋与明遗民网络，把海上基地转化为更大的反清战场。围城并不只考验冲锋；军队要在内河环境中维持粮食、船只、城外营地和联络，清军则集中江防、城门与守备力量。郑军未能突破城防而撤回，显示其攻势在内河攻坚与长期补给面前难以持续。清廷记录、杨英《从征实录》与荷兰海上报告都能确认北伐、围城和撤退的主干，但内应的范围、作战计划及伤亡互有歧异；把这次行动写成几乎成功的复明，或写成毫无准备的冒险，都会超出材料所能承担的程度。

撤退没有使台湾自动成为唯一的去处，却使那里成为更迫切的战略选择。顺治十八年（1661年），清廷命福建、广东沿海迁界，意在把郑氏同岸上的粮食、人口与信息来源隔开。谕令所能直接告诉我们的，是中央选择了以强制迁徙制造隔离带；它不能替我们证明每一县迁了多远、多少人死于途中，或秘密贸易有多大规模。沿海县志和通志留下的迁撤、复界线索，使渔盐、耕作、宗族与港口关系被打断的代价可以进入叙事；但这些记载的密度和视角各异，不能把局部遭遇拼成一幅均一的海岸图景。

迁界也并未把海面与陆地彻底分开。禁令增加了往来成本，海岸居民、地方官和海上经营者仍可能在不同地点以不同方式应对：有人被迫内迁，有人守着可耕之地或亲族关系，有人继续依赖港口和航路。现有材料足以说明隔离带未能完全阻绝往来，却不足以量化这些变通，更不足以把它们叙述成全海岸一致、持续不断的走私网络。政策的直接军事目的，是压缩郑氏补给的可能；它的副作用，则是把本已被战争牵动的生计、迁徙和地方关系再度推入不确定之中。

同年四月，郑成功军登陆台湾，围攻荷兰东印度公司在热兰遮城的据点。对郑氏而言，这并非单纯向外扩张：南京行动受挫之后，需要一处能够安置军政、取得粮食并控制航路的基地；对岛上的荷兰据点、汉人移民社会与原住民社群而言，战争则把既有的土地、征发与交往关系重新置于强制之下。围城的胜负取决于守军补给能否进入、周边村社与航道由谁控制，以及攻守双方能否把人力和粮食留在城下，而不只是城墙前的一次决战。

次年正月，热兰遮城守军投降，荷兰东印度公司在台湾西南的统治至此终结。投降条约和《热兰遮城日志》能够确定交接的时点，也能让人看见围城中荷兰方面的处境；清廷军报与后出的史书则呈现北京获知这一消息的路径。它们不能替城外居民、被俘者或原住民社群作证。兵员、联盟和伤亡在中荷战报中各有立场，因而不能用任一方的数字填满这场转换。可以确认的是，郑氏接收了城堡、田地和贸易节点，台湾由此成为反清政权筹集人口与财政资源的基地；至于这种接收怎样抵达岛上不同社群、各处又怎样回应，仍须依赖保存不均的地方、契约和多语材料逐地重建。

郑成功在康熙元年五月去世，使新得的基地立刻面对继承问题。城堡已经易手，并不表示军政、宗族与海商网络已经形成稳定的服从关系；郑经最终掌权，是在这些网络重新安排权力与资源的过程中出现的结果。有关死因、继承争执和军中处置的细节，主要依赖后出的郑氏材料与荷兰报告，叙事只能把握重组的压力，不能替当事人补出一致的动机。

康熙二年（1663年），清军联合部分郑氏降将攻取厦门、金门，郑氏主力退守台湾。熟悉海上条件的降将与沿海据点帮助清军夺取近大陆的前进基地，但厦金易手并不等于清廷已经取得制海权。对郑氏而言，失去两岛迫使台湾的土地、港口和跨海贸易承担更重的军政负担；对清廷而言，迁界所制造的隔离、沿海守备与水师建设仍要持续付出社会与后勤成本。这里的制度得失并不在于哪一方终于“控制了海岸”，而在于双方都不得不以迁徙、城防、航路和粮饷来维持一场尚未结束的海上战争。

